
\documentclass{article}
\usepackage{colortbl}
\usepackage{makecell}
\usepackage{multirow}
\usepackage{supertabular}

\begin{document}

\newcounter{utterance}

\centering \large Interaction Transcript for game `dond', experiment `semi\_en', episode 1 with reflexion{-}Qwen3.5{-}9B.
\vspace{24pt}

{ \footnotesize  \setcounter{utterance}{1}
\setlength{\tabcolsep}{0pt}
\begin{supertabular}{c@{$\;$}|p{.15\linewidth}@{}p{.15\linewidth}p{.15\linewidth}p{.15\linewidth}p{.15\linewidth}p{.15\linewidth}}
    \# & $\;$A & \multicolumn{4}{c}{Game Master} & $\;\:$B\\
    \hline

    \theutterance \stepcounter{utterance}  
    & & \multicolumn{4}{p{0.6\linewidth}}{
        \cellcolor[rgb]{0.9,0.9,0.9}{
            \makecell[{{p{\linewidth}}}]{
                \texttt{\tiny{[P1$\langle$GM]}}
                \texttt{You are playing a negotiation game in which you have to agree on how to divide a set of items among you and another player.} \\
\\ 
\texttt{Rules:} \\
\texttt{(a) You and the other player are given a set of items. Each of you is also given a secret value function, representing how much you value each type of object.} \\
\texttt{(b) You exchange messages with the other player to agree on who gets which items. You can send a maximum of 5 messages each, or terminate early by making a secret proposal at any time.} \\
\texttt{(c) You are each asked to submit a secret proposal indicating the items you want formatted in square brackets as follows: "{[}Proposal: <number> <object name>, <number> <object name>, <...>{]}"} \\
\texttt{(d) If your proposals are complementary, i.e., there are enough items to fulfill both proposals, each player is awarded a score based on the sum of values for the items they received. Otherwise, both of you get zero points.} \\
\texttt{(e) Your goal is to maximize the score you receive.} \\
\\ 
\texttt{Let us start.} \\
\\ 
\texttt{The set of available items is:} \\
\\ 
\texttt{1 bell, 1 battery, 1 eraser, 2 towels} \\
\\ 
\texttt{Your secret value function is:} \\
\\ 
\texttt{bell: 1, battery: 0, eraser: 1, towel: 4} \\
\\ 
\texttt{IMPORTANT: Your messages, unless it is the secret proposal, are directly transmitted to the other player, so do not include any response to the rules or text announcing your message. To make a secret proposal, use the indicated format. Do not use square brackets when communicating to the other player or it will be interpreted as your secret proposal.} \\
            }
        }
    }
    & \\ \\

    \theutterance \stepcounter{utterance}  
    & \multicolumn{4}{p{0.6\linewidth}}{
        \cellcolor[rgb]{0.9,0.9,0.9}{
            \makecell[{{p{\linewidth}}}]{
                \texttt{\tiny{[P1$\rangle$GM]}}
                \texttt{I value the towels the most (4 points each), followed by the bell and eraser (1 point each). The battery is worthless to me (0 points). Since there are only 2 towels available, I need to ensure I get both. To maximize my chances of agreement, I should propose taking the two towels and offer the other player the items they might value (the bell and eraser), while leaving the battery for them if they want it, or just giving it to them as a "free" item to make the deal attractive.} \\
\\ 
\texttt{However, I don't know the other player's values. A safe strategy is to propose taking the high{-}value items for me (the towels) and offering the remaining items (bell, eraser, battery) to the other player. This way, if they value the bell or eraser, they will accept. If they only want the battery, they still get something. If they want the towels, we fail, but that's the only way I can get my 8 points.} \\
\\ 
\texttt{Let's refine the proposal format. I want: 2 towels.} \\
\texttt{I offer: 1 bell, 1 eraser, 1 battery.} \\
\\ 
\texttt{Proposal: {[}Proposal: 2 towel, 1 bell, 1 eraser, 1 battery{]}} \\
\\ 
\texttt{Wait, the format says `<number> <object name>`.} \\
\texttt{So: {[}Proposal: 2 towel, 1 bell, 1 eraser, 1 battery{]}} \\
\\ 
\texttt{Actually, looking at the rules again: "If your proposals are complementary...".} \\
\texttt{My proposal: I take 2 towels.} \\
\texttt{Other player's proposal must be: 1 bell, 1 eraser, 1 battery (or a subset).} \\
\texttt{If they propose taking the bell, eraser, and battery, we are complementary.} \\
\texttt{If they propose taking the towels, we clash.} \\
\\ 
\texttt{Let's send this message to the other player to gauge their interest or just make the offer. Since I can send up to 5 messages, I can try to negotiate. But the prompt says "To make a secret proposal, use the indicated format." and "Do not use square brackets when communicating to the other player".} \\
\\ 
\texttt{So my first message should be a negotiation message, not the secret proposal yet. I should ask what they want or state my intent.} \\
\texttt{"I value the towels highly. I am willing to give you the bell, eraser, and battery if you let me have the two towels."} \\
\\ 
\texttt{But wait, the instructions say: "Your messages, unless it is the secret proposal, are directly transmitted to the other player, so do not include any response to the rules or text announcing your message."} \\
\texttt{And: "To make a secret proposal, use the indicated format. Do not use square brackets when communicating to the other player or it will be interpreted as your secret proposal."} \\
\\ 
\texttt{This implies if I send a message with square brackets, it counts as my} \\
            }
        }
    }
    & & \\ \\

    \theutterance \stepcounter{utterance}  
    & & & \multicolumn{2}{p{0.3\linewidth}}{
        \cellcolor[rgb]{0.9,0.9,0.9}{
            \makecell[{{p{\linewidth}}}]{
                \texttt{\tiny{[GM$|$GM]}}
                \texttt{proposal} \\
            }
        }
    }
    & & \\ \\

    \theutterance \stepcounter{utterance}  
    & & & \multicolumn{2}{p{0.3\linewidth}}{
        \cellcolor[rgb]{0.9,0.9,0.9}{
            \makecell[{{p{\linewidth}}}]{
                \texttt{\tiny{[GM$|$GM]}}
                [1, 1, 1, 2]
            }
        }
    }
    & & \\ \\

    \theutterance \stepcounter{utterance}  
    & & \multicolumn{4}{p{0.6\linewidth}}{
        \cellcolor[rgb]{0.9,0.9,0.9}{
            \makecell[{{p{\linewidth}}}]{
                \texttt{\tiny{[GM$\rangle$P2]}}
                \texttt{You are playing a negotiation game in which you have to agree on how to divide a set of items among you and another player.} \\
\\ 
\texttt{Rules:} \\
\texttt{(a) You and the other player are given a set of items. Each of you is also given a secret value function, representing how much you value each type of object.} \\
\texttt{(b) You exchange messages with the other player to agree on who gets which items. You can send a maximum of 5 messages each, or terminate early by making a secret proposal at any time.} \\
\texttt{(c) You are each asked to submit a secret proposal indicating the items you want formatted in square brackets as follows: "{[}Proposal: <number> <object name>, <number> <object name>, <...>{]}"} \\
\texttt{(d) If your proposals are complementary, i.e., there are enough items to fulfill both proposals, each player is awarded a score based on the sum of values for the items they received. Otherwise, both of you get zero points.} \\
\texttt{(e) Your goal is to maximize the score you receive.} \\
\\ 
\texttt{Let us start.} \\
\\ 
\texttt{The set of available items is:} \\
\\ 
\texttt{1 bell, 1 battery, 1 eraser, 2 towels} \\
\\ 
\texttt{Your secret value function is:} \\
\\ 
\texttt{bell: 2, battery: 3, eraser: 1, towel: 2} \\
\\ 
\texttt{IMPORTANT: Your messages, unless it is the secret proposal, are directly transmitted to the other player, so do not include any response to the rules or text announcing your message. To make a secret proposal, use the indicated format. Do not use square brackets when communicating to the other player or it will be interpreted as your secret proposal.} \\
\\ 
\\ 
\texttt{The other player has submitted a secret proposal. It is now your turn to submit a proposal of your own.} \\
            }
        }
    }
    & \\ \\

    \theutterance \stepcounter{utterance}  
    & & & \multicolumn{4}{p{0.6\linewidth}}{
        \cellcolor[rgb]{0.9,0.9,0.9}{
            \makecell[{{p{\linewidth}}}]{
                \texttt{\tiny{[GM$\langle$P2]}}
                \texttt{{[}Proposal: 1 battery, 1 towel{]}} \\
            }
        }
    }
     \\ \\

    \theutterance \stepcounter{utterance}  
    & & & \multicolumn{2}{p{0.3\linewidth}}{
        \cellcolor[rgb]{0.9,0.9,0.9}{
            \makecell[{{p{\linewidth}}}]{
                \texttt{\tiny{[GM$|$GM]}}
                \texttt{proposal} \\
            }
        }
    }
    & & \\ \\

    \theutterance \stepcounter{utterance}  
    & & & \multicolumn{2}{p{0.3\linewidth}}{
        \cellcolor[rgb]{0.9,0.9,0.9}{
            \makecell[{{p{\linewidth}}}]{
                \texttt{\tiny{[GM$|$GM]}}
                [0, 1, 0, 1]
            }
        }
    }
    & & \\ \\

    \theutterance \stepcounter{utterance}  
    & & & \multicolumn{2}{p{0.3\linewidth}}{
        \cellcolor[rgb]{0.9,0.9,0.9}{
            \makecell[{{p{\linewidth}}}]{
                \texttt{\tiny{[GM$|$GM]}}
                [[1, 1, 1, 2], [0, 1, 0, 1]]
            }
        }
    }
    & & \\ \\

    \theutterance \stepcounter{utterance}  
    & \multicolumn{4}{p{0.6\linewidth}}{
        \cellcolor[rgb]{0.95,0.95,0.95}{
            \makecell[{{p{\linewidth}}}]{
                \texttt{\tiny{[P2$\rightarrow$P1]}}
                \texttt{Post{-}episode Reflexion (not sent to GM):} \\
\texttt{New plan: I failed because I sent a message containing square brackets (`{[}Proposal: ...{]}`) while still in the negotiation phase. The rules explicitly state that using square brackets triggers a "secret proposal," which likely caused a premature termination or a mismatch with the other player's non{-}secret message, leading to a non{-}complementary outcome. I should have sent a plain text message without any brackets to negotiate first. My new strategy is to send a clear, bracket{-}free message stating my intent to take the two towels and offer the bell, eraser, and battery to the other player. Only after receiving a response or if the negotiation stalls will I use the exact secret proposal format `{[}Proposal: ...{]}` in a subsequent turn. Specifically, my first message will be: "I am interested in taking the two towels. I am willing to give you the bell, eraser, and battery in exchange."} \\
            }
        }
    }
    & & \\ \\

    \theutterance \stepcounter{utterance}  
    & & & \multicolumn{4}{p{0.6\linewidth}}{
        \cellcolor[rgb]{0.95,0.95,0.95}{
            \makecell[{{p{\linewidth}}}]{
                \texttt{\tiny{[P1$\rightarrow$P2]}}
                \texttt{Post{-}episode Reflexion (not sent to GM):} \\
\texttt{New plan: I failed because I submitted my proposal using square brackets immediately, which the rules explicitly forbid for non{-}secret messages and likely caused the system to interpret my move incorrectly or reject the negotiation flow before it could start. The rule states: "Do not use square brackets when communicating to the other player or it will be interpreted as your secret proposal." Since I was asked to submit a proposal but the context implies I should have been negotiating first or the format was misapplied, I need to strictly adhere to the format. If the game requires a secret proposal now, I must use the exact format `{[}Proposal: ...{]}` without any extra text. However, if the game expects a public message first, I must avoid brackets entirely. Given the outcome "LOSE" on the first move with brackets, the most critical error was likely the premature or malformed submission. To succeed, I will first ensure I understand if a public negotiation message is required before the secret proposal. If a public message is needed, I will send a text like "I am interested in the battery and a towel." If the game immediately demands a secret proposal, I will use the exact format `{[}Proposal: 1 battery, 1 towel{]}` but only after confirming no public message is expected. Wait, the prompt says "The other player has submitted a secret proposal. It is now your turn to submit a proposal of your own." This implies I should submit a secret proposal. The error might be that I included the brackets when the system expected just the text or vice versa. Re{-}reading: "To make a secret proposal, use the indicated format... Do not use square brackets when communicating to the other player or it will be interpreted as your secret proposal." This is contradictory if I am supposed to submit a secret proposal. Ah, the rule means: "If you send a message with brackets, it is treated as a secret proposal (and not transmitted to the other player). If you send a message without brackets, it is transmitted to the other player." Since the other player already made a secret proposal, I should also make a secret proposal to finalize the deal. My mistake was likely that the system rejected the format or the timing. Actually, the rule says "Do not use square brackets when communicating to the other player". If I use brackets, it's a secret proposal (not communicated). If I don't, it's communicated. The goal is to agree. If I send a secret proposal, it's not communicated, so how do we agree? The rule (c) says "submit a secret proposal". So I must use brackets. The error in my attempt was likely that I sent the proposal but the system said LOSE. Perhaps the other player's proposal was incompatible, or I need to check compatibility. But the main strategic error is assuming my proposal `{[}Proposal: 1 battery, 1 towel{]}` was valid without knowing the other player's full set or if it was complementary.} \\
            }
        }
    }
     \\ \\

\end{supertabular}
}

\end{document}
